\noindent\textit{2015 manuscript.}

\section{Joule-Thomson Effect}

cf. Problem 1.7.6 of Le Bellac, Mortessagne, Batrouni (2004) \cite{MLeBellacFMortessagneGBatrouni2004}, Lecture 18, Sections ``Throttling Processes'', ``The Joule-Thomson Effect'' of Groth, Phys301 Fall 1999 \footnote{Edward J. Groth, 1999, Fall 1999 Physics 301, \href{http://grothserver.princeton.edu/~groth/phys301f99/lect18.pdf}{Phys301 23-Nov-1999 18-10} }

From Prob. 1.7.6 of Le Bellac, Mortessagne, Batrouni (2004) \cite{MLeBellacFMortessagneGBatrouni2004}, their setup is the following: \\
gas initially confined in fixed volume escapes freely without exchanging heat with its environment, 
\begin{itemize}
\item initially, gas in compartment and $(T_1,V_1,p_1)$ 
\item expands into container of volume $V_2$ under constant pressure $p_2$ ($p_2<p_1$)
\item 2 compartments are thermally isolated; connected by porous plug to maintain pressure difference
\item adiabatic (but not quasi-static) expansion.
\end{itemize}

$Q=0$.  So $Q=0 = dU -W$.  Then
\[
\int dU - \int W = \Delta U - \int W = 0 
\]

Now
\[
\Delta H = H_2 - H_1 = u_2 + p_2 V_2 - U_1 - p_1 V_1 = \Delta U + p_2 V_2 - p_1 V_1 
\]
Treat the process of gas evacuating, or going from compartment 1 to container 2 as 2 processes:  $\int W = W_1 + W_2$ where \\
$W_1 = -p_1 (0-V_1) = p_1 V_1 \equiv $ work done on gas as the gas evacuates the compartment. \\
$W_2 = -p_2(V_2-0) = -p_2 V_2 = $ work done by gas as gas expands into container at constant pressure $p_2$.  

Thus $\int W = W_1 + W_2 = p_1 V_1 - p_2 V_2$ and so 
\[
\boxed{ \Delta H = 0 }
\]
The enthalpy is constant for this expansion.  

For an ideal gas, $H= U + pV = \frac{\gamma}{\gamma -1} N\tau$, for an ideal gas has $U = \frac{1}{\gamma -1} N\tau$.  $\gamma = \frac{5}{3}$ and $\frac{\gamma}{\gamma-1} = \frac{5/3}{ 2/3} = 5/2$.  Nevertheless, since $H$ is constant, $T$ is constant, and so there's no cooling of an ideal gas in a Joule-Thomson apparatus!  But an ideal gas has no interaction energy, and it can be argued that one has to consider interaction energy between molecules for liquefaction (cf. Groth (1999) \cite{EGroth1999}, \href{http://grothserver.princeton.edu/~groth/phys301f99/lect18.pdf}{Lecture 18 Physics 301 Fall 1999}).  

From the ``Derivation of the Joule-Thomson (Kelvin) coefficient'' section of the \emph{Joule-Thomson effect} article in \href{https://en.wikipedia.org/wiki/Joule–Thomson_effect}{Wikipedia}, recall the definition of enthalpy $H:= U + pV$, and so that $dH= dU+ pdV + Vdp=\tau d\sigma + Vdp$.  

Let $\sigma = \sigma(\tau,p)$.  
\[
\begin{gathered}
  dH = \tau \left( \frac{ \partial \sigma}{ \partial \tau } \right)_p d\tau + (V + \tau \left( \frac{ \partial \sigma }{ \partial p }\right)_{\tau} ) dp = C_p d\tau + (V + \tau \left( \frac{ \partial \sigma }{ \partial p }\right)_{\tau} ) dp
\end{gathered}
\]
since
\[
\begin{gathered}
Q = Q(\tau,p) = \frac{ \partial Q}{ \partial \tau} d\tau + \frac{ \partial Q}{ \partial p } dp = C_p d\tau + \frac{ \partial Q}{ \partial p} dp
\end{gathered}
\]
and so for a thermodynamic process $\gamma$ such that $dp(\dot{\gamma})=0$ (i.e. constant pressure), then we can say
\[
\frac{C_p}{\tau}d\tau = d\sigma \xrightarrow{ \int_{\gamma } } \int_{\gamma} \frac{C_p}{\tau} d\tau = \Delta \sigma \Longrightarrow \tau \left( \frac{ \partial \sigma }{ \partial \tau } \right)_p = C_p
\]
And now 
\[
dG = -\sigma d\tau  + Vdp = \left( \frac{ \partial G}{ \partial \tau } \right)_p d\tau + \left( \frac{ \partial G}{ \partial p } \right)_{\tau} dp
\]
and since $\frac{ \partial^2 G}{ \partial \tau \partial p} = \frac{ \partial^2 G}{ \partial p \partial \tau}$, 
\[
\left( \frac{ \partial \sigma }{ \partial p } \right)_{\tau} = - \left( \frac{ \partial V}{ \partial \tau} \right)_p = -V\alpha
\]
Then
\[
dH = C_p d\tau + V(1-\tau \alpha) dp
\]
and so 
\begin{equation}\label{Eq:JouleThomsoncoeff}
\left( \frac{ \partial \tau}{ \partial p }\right)_H = \frac{ V(1-\tau \alpha )}{-C_p}
\end{equation}
is the \textbf{Joule-Thomson coefficient}, where $\alpha$ is the coefficient of thermal expansion.  

For an ideal gas, 
\[
\begin{gathered}
  pV = N\tau \\ 
  V = \frac{N\tau}{p}
\end{gathered} \Longrightarrow \frac{ \partial V}{ \partial \tau} = \frac{N}{p} = \frac{V}{\tau}
\]
so
\[
dH = C_pd\tau + ( V + \tau \left( \frac{-V}{\tau} \right) ) dp = C_p d\tau 
\]
so that 
\[
\left( \frac{ \partial \tau}{ \partial p} \right)_H = 0 
\]
for an ideal gas.  

For every real gas, every real gas has an \emph{inversion temperature} such that \\
below the inversion temperature, $\left( \frac{ \partial T}{ \partial p} \right)_H$ is positive, and so $dp <0$ always in an expansion, and so $dT <0$ and the gas cools, and \\
above the inversion temperature, $\left( \frac{ \partial T}{ \partial p} \right)_H$ is negative, and so $dp <0$ always in an expansion, and so $dT <0$ and the gas warms.
\begin{correction}{Joule--Thomson warming}
Above inversion the coefficient is negative and $dp<0$, so $d\tau>0$.
The gas warms.
The clause $dT<0$ in the sentence has the wrong sign.
See \S\ref{sec:jt-sign}.
\end{correction}  

An inversion temperature must exist because when the molecules are close enough, there are intermolecular forces that'll make them attract each other (such as temporary configurations so a slightly (electrically) negative cloud of electrons attract slightly positively charged part of another molecule) so the interaction energy is negative.  

However, it appears that this inversion temperature is not found from the Joule-Thomson coefficient expression, Eq. \ref{Eq:JouleThomsoncoeff}, but from, possibly, the Van der Waals expression or another model (cf. Lecture 18 of Groth (1999) \cite{EGroth1999}).  

Looking at Figure 1.16 of  Le Bellac, Mortessagne, Batrouni (2004) \cite{MLeBellacFMortessagneGBatrouni2004}, for isenthalpy curves ($H=$ const.), in $(T,p)$ plane, $T$ vs. $p$, the dashed line intersects the points on each isenthalpy curves where $\left( \frac{ \partial T}{ \partial p } \right)_H =0$ and for lower pressures ($dp<0$, gas expands), cooling occurs.  

There's a maximum temperature, about $300^{\circ} \, C$ for nitrogen, where the dashed line meets the $p=0 \, \text{atm}$ axis, meaning $\left( \frac{ \partial \tau}{ \partial p} \right)_H <0$ \, $\forall \, p >0$, so no cooling occurs.  

